\section{环境保护法}
\begin{enumerate}
  \item 年头久。
  \item 生态环境主管部门 \yjt 监管部门。质量负责 \yjt 各级政府。
  \item 基本制度。
        \begin{enumerate}
          \item 3 同时，建设项目中的防污设施应当与主体工程同时设计、施工、使用。不得擅自拆除、闲置。
          \item 环境。
                \begin{enumerate}
                  \item 环境保护税。直接向环境排放、超标排放。
                  \item 环境质量标准。水、大气、土壤。容许浓度。
                        \begin{enumerate}
                          \item 是制定污染物排放标准的依据。
                          \item 分类和制定。\yi 国标 \yjt 国务院生态环境主管部门 。\er 地方标准 \yjt 省政府 \yjt 国务院备案。国无我有，国有我严。
                        \end{enumerate}
                  \item 农业环境保护。不管它。
                \end{enumerate}
          \item 污染。
                \begin{enumerate}
                  \item 排污许可管理。
                  \item 排污收费。达不达标都要交。（环境保护）税、费不并行。
                  \item 污染物排放标准。人为污染源。\yi 是实现环境质量标准的重要保证。\er 分类和制定。同上。
                  \item 重点污染物排放总量控制。
                        \begin{enumerate}
                          \item 3 级拆解。国务院下达 \yjt 省政府分解落实 \yjt 企业。
                          \item 超过总量控制目标、未完成环境质量目标 \yjt 应当暂停审批，不得开工建设。
                        \end{enumerate}
                \end{enumerate}
          \item 生态。
                \begin{enumerate}
                  \item 生态红线。重点生态功能区、生态环境敏感区、脆弱区。
                  \item 保护生态多样性。外来物种引进，采取措施，防止破坏。
                  \item 生态保护补偿机制。跨省；加大转移支付；国家协商，受益地区和生态保护地区政府的补偿机制。
                \end{enumerate}
          \item 信息公开。国重大、省公报、县违法、重自己、见公征、审环评。
                \begin{enumerate}
                  \item 国务院 \yjt 重大。
                  \item 省 \yjt 公报、本省整体情况。
                  \item 县以上 \yjt 违法者名单。
                  \item 重点排污单位 \yjt 自己。
                  \item 建设单位 \yjt 公众、征求意见、说明情况。
                  \item 审批环评部门 \yjt 环评文件。
                \end{enumerate}
        \end{enumerate}
  \item 环境民事责任。同民法。
  \item 环境民事公益诉讼。设区的市以上、民政部门登记的、社会组织 + 5 年以上没有违法记录。
  \item 环境行政责任。当行政法题目去做。常见考点。
        \begin{enumerate}
          \item 行政处罚。
                \begin{enumerate}
                  \item 县以上环保部门职权。行为罚。现场检查；查扣；限制生产；责令停业、停产、关闭（政府批）。
                  \item 按日连续处罚 3 步走。小重点。
                        \begin{enumerate}
                          \item 违法排污 \yjt 受到了罚款处罚 \yjt 责令改正拒不改正。责令改正次日起算。
                          \item 地方性法规可以增设情形（\ydy 新花样）。
                        \end{enumerate}
                \end{enumerate}
          \item 行政措施。\yi 移送公安。\er 主管人员、直接责任人员，公安拘留。
                \begin{enumerate}
                  \item 轻的 \yjt 拒不改正 \yjt 行政措施。\al{未经过环评开工、未取得排污许可证、使用禁止的农药}。
                  \item 重的 \yjt 直接行政措施。\al{逃避监管、造假}。
                \end{enumerate}
        \end{enumerate}
\end{enumerate}
